\documentclass[12pt]{article}
\usepackage[T1]{fontenc}
\usepackage[utf8]{inputenc}
\usepackage{lmodern}
\usepackage{textcomp}
\usepackage{enumitem}
\usepackage{geometry}
\geometry{margin=1in}
\begin{document}
\begin{center}
Answer Key - Constitutional Law - Undergraduate Level Assessment 2 \
\small (Answers are ordered exactly as in the question paper.)
\end{center}

\section*{Multiple Choice}
\begin{enumerate}[label=\arabic*.]
\item \textbf{Answer:} C
\\ \textit{Options:} A) There is no relationship between law and morality.; B) Law is always in advance of moral ideas.; C) The law is inextricably bound up with morals.; D) Morality is generally in advance of the law.
\\ \textit{Note:} Correct option: C - The law is inextricably bound up with morals.
\item \textbf{Answer:} B
\\ \textit{Options:} A) Because its justification depends on the concept employed.; B) Because any definition of punishment should be value-neutral.; C) Because the concept of punishment is controversial.; D) Because the practice of punishment is separate from its justification.
\\ \textit{Note:} Correct option: B - Because any definition of punishment should be value-neutral.
\item \textbf{Answer:} C
\\ \textit{Options:} A) In hard cases judges generally decide cases on the basis of rights.; B) The rights of the parties feature in the determination of most cases before the courts.; C) Judges exercise strong discretion.; D) Judges seek the best 'fit' with constitutional and institutional history.
\\ \textit{Note:} Correct option: C - Judges exercise strong discretion.
\item \textbf{Answer:} A
\\ \textit{Options:} A) The POP will choose wealth over a compassionate society.; B) The POP will choose equality over power.; C) The POP will be unselfish.; D) The POP will choose to protect the disabled.
\\ \textit{Note:} Correct option: A - The POP will choose wealth over a compassionate society.
\item \textbf{Answer:} D
\\ \textit{Options:} A) Incorporationism' accepts that judges may decide cases by reference to moral factors.'; B) A soft positivist rejects the role of morality in the description of law.; C) Sometimes the definition of law includes moral considerations.; D) Judges lack strong discretion.
\\ \textit{Note:} Correct option: D - Judges lack strong discretion.
\end{enumerate}

\section*{True / False}
\begin{enumerate}[label=\arabic*.]
\setcounter{enumi}{5}
\item \textbf{Answer:} False
\\ \textit{Options:} A) True; B) False
\\ \textit{Note:} LLM-generated false statement. Correct answer: 'Formalism treats legal reasoning as syllogistic reasoning.'.
\item \textbf{Answer:} True
\\ \textit{Options:} A) True; B) False
\\ \textit{Note:} LLM-generated true statement based on MCQ.
\item \textbf{Answer:} False
\\ \textit{Options:} A) True; B) False
\\ \textit{Note:} LLM-generated false statement. Correct answer: 'Purely discussed jurisprudence only'.
\item \textbf{Answer:} False
\\ \textit{Options:} A) True; B) False
\\ \textit{Note:} LLM-generated false statement. Correct answer: 'Analytical School'.
\item \textbf{Answer:} False
\\ \textit{Options:} A) True; B) False
\\ \textit{Note:} LLM-generated false statement. Correct answer: 'If each person's holdings are just, then the total distribution of holdings is just.'.
\end{enumerate}

\section*{Long Form Response}
\begin{enumerate}[label=\arabic*.]
\setcounter{enumi}{10}
\item \textbf{Answer:} To provide housing for all those homeless within a state. Explain why this is the correct answer and provide supporting evidence. Reference key facts or concepts that support this choice.
\\ \textit{Note:} Long-form derived from MCQ; award credit for stating the correct idea and giving supporting reasoning.
\item \textbf{Answer:} Mediation is usually conducted by a person appointed with the consent of the parties, while conciliation involves a commission, which proceeds to an impartial examination of the dispute and proposes settlement terms. Explain why this is the correct answer and provide supporting evidence. Reference key facts or concepts that support this choice.
\\ \textit{Note:} Long-form derived from MCQ; award credit for stating the correct idea and giving supporting reasoning.
\end{enumerate}

\vfill
\noindent\textit{End of Answer Key}
\end{document}